\documentclass[11pt]{article}
\usepackage[T1]{fontenc}
\usepackage[utf8]{inputenc}
\usepackage{lmodern}
\usepackage{microtype}
\usepackage[margin=1in]{geometry}
\usepackage{amsmath,amssymb,amsthm,mathtools}
\usepackage{booktabs,array,longtable}
\usepackage{xcolor}
\usepackage{enumitem}
\usepackage{fancyhdr}
\usepackage{hyperref}
\definecolor{inkblue}{RGB}{28,54,78}
\hypersetup{colorlinks=true,linkcolor=inkblue,citecolor=inkblue,urlcolor=inkblue,
 pdftitle={The Cubic Boundary for Restricted Partitions},
 pdfauthor={Research report prepared for Vladimir Reshetnikov}}
\setlength{\headheight}{14pt}
\pagestyle{fancy}
\fancyhf{}
\fancyhead[L]{\small\textsc{The cubic boundary for restricted partitions}}
\fancyhead[R]{\small Research report}
\fancyfoot[C]{\thepage}
\setlength{\parskip}{4pt}
\setlength{\parindent}{1.15em}
\setlist{itemsep=3pt,topsep=5pt}
\newtheorem{theorem}{Theorem}[section]
\newtheorem{lemma}[theorem]{Lemma}
\newtheorem{proposition}[theorem]{Proposition}
\newtheorem{corollary}[theorem]{Corollary}
\theoremstyle{definition}
\newtheorem{definition}[theorem]{Definition}
\theoremstyle{remark}
\newtheorem{remark}[theorem]{Remark}
\newcommand{\fall}[2]{#1^{\underline{#2}}}
\newcommand{\E}{\mathbb E}
\newcommand{\Prob}{\mathbb P}
\newcommand{\Z}{\mathbb Z}
\newcommand{\N}{\mathbb N}
\newcommand{\ii}{\mathrm i}
\newcommand{\dd}{\,\mathrm d}
\DeclareMathOperator{\Pois}{Poisson}
\DeclareMathOperator{\Fix}{Fix}
\DeclareMathOperator{\Var}{Var}
\numberwithin{equation}{section}
\newenvironment{writenote}[1][]{\par\smallskip\begingroup\small
 \noindent\textbf{[write]}\if\relax\detokenize{#1}\relax\else~\textbf{#1.}\fi\ }%
 {\par\endgroup\smallskip}

\title{\color{inkblue}\textbf{The Cubic Boundary for Restricted Partitions}\\[5pt]
\large Sharp volume laws, uniform all-order expansions,\\
exponentially small corrections, inverse asymptotics,\\
and a Poisson law for repeated parts}
\author{Research report prepared for Vladimir Reshetnikov}
\date{1 October 2026}

\begin{document}
\maketitle
\begin{abstract}
Let $p_m(N)$ count partitions of $N$ into parts at most $m$, equivalently
partitions into at most $m$ parts. We give an elementary necessary and sufficient
condition for the volume approximation:
\[
 p_m(N)\sim\frac{N^{m-1}}{m!(m-1)!}
 \quad\Longleftrightarrow\quad \frac{N}{m^3}\longrightarrow\infty,
 \qquad m\longrightarrow\infty.
\]
This proves, and sharpens, a condition still described as conjectural in
OEIS A238016 and A258670. The sufficiency has classical antecedents and is not
claimed as a new discovery. At the critical scale $N\sim c m^3$, the ratio to
volume tends to $e^{1/(4c)}$. We prove a centered expansion uniform for
$N\ge c_0m^3$, give an all-order coefficient algorithm and five explicit
corrections for A238608, and derive a second expansion that resolves actual
exponential sectors for $N=q^m$. Both types are inverted, with careful
separation of smooth asymptotic charts from integer-valued inverse functions.
In the conjugate model with at most $m$ parts, the number of equal adjacent
parts converges to $\Pois(1/(2c))$; its probability generating function has an
explicit first finite-size correction. Exact computation checks 258 samples,
including all displayed initial A238608 values, and the package supplies
symbolic derivation and numerical verification code.
\end{abstract}

\noindent\textbf{Status.} This is a self-contained proof-and-extension report.
Its leading partition formulas and first-wave machinery have substantial
classical precedents. The explicit refinements below are derived here, but no
exhaustive priority claim is made. The proofs have not been independently
peer reviewed or formalized in Lean. Exact finite checks corroborate formulas;
they do not replace the uniform estimates proved in the text.

\clearpage
\begingroup
\setlength{\parskip}{0pt}
\tableofcontents
\endgroup
\clearpage

\section{The OEIS problem, the source audit, and the scope of the results}\label{rpc:sec:intro}

\subsection{One fundamental family, rather than unrelated sequences}

Throughout, $m$ and $N$ are positive integers and $m\to\infty$. Define
\begin{equation}\label{rpc:eq:gf}
 F_m(q)=\prod_{j=1}^m\frac1{1-q^j},\qquad
 p_m(N)=[q^N]F_m(q),\qquad
 V_m(N)=\frac{N^{m-1}}{m!(m-1)!}.
\end{equation}
The equivalence between the two partition interpretations follows by
conjugating Ferrers diagrams. This is a standard, widely useful family of
restricted partitions, not an isolated artificial recurrence. The same
analytic problem occurs when the size $N$ grows polynomially, exponentially,
or factorially with the allowed number $m$ of parts.

The following are the specific OEIS targets inspected for this report.
The descriptions and posted leading formulas are recorded in the respective
entries \cite{oeis238016,oeis238608,oeis238010,oeis258670}.
\begin{center}
\small
\begin{tabular}{@{}p{0.23\textwidth}p{0.28\textwidth}p{0.41\textwidth}@{}}
\toprule
Entry or family & Count & Role in this report\\
\midrule
A238016 & $p_m(m^k)$, an array & A posted general conjecture; cubic threshold\\
A238608 & $p_m(m^3)$ & Five corrections and inverse expansion\\
A258302--A258305 & $p_m(c m^3)$, $c=2,3,4,5$ & Uniform critical-scale family\\
A238609--A238615 & $p_m(m^k)$, $4\le k\le10$ & All-order power-column expansions\\
A238010; A237998 & $p_m(q^m)$; $q=2$ & Exponential sectors and their inversion\\
A238000 & $p_m(m^m)$ & Superexponential sectors and a Lambert chart\\
A258670 & $p_m((2m)!)$ & Factorial-scale consequence\\
\bottomrule
\end{tabular}
\end{center}

The general condition in A238016 is expressed using the nonstandard notation
$f(m)\ge O(m^4)$. Interpreting the intended assertion as growth at least a
positive constant times $m^4$, it asserts $p_m(f(m))\sim V_m(f(m))$.
The exact criterion in Theorem~\ref{rpc:thm:iff} is instead $f(m)/m^3\to\infty$.
It covers, for example, $f(m)=\lceil m^3\log m\rceil$, not merely quartic growth.

\subsection{What is classical, and what this report actually adds}

A condition remaining conjectural on an OEIS page does not establish that
it is an open problem in the research literature. Canfield's account of
Szekeres' formula records the Erd\H{o}s--Lehner estimate
\begin{equation}\label{rpc:eq:classical}
 p_m(N)=\frac1{m!}\binom{N-1}{m-1}
          \exp\!\left(O(m^3/N)\right),\qquad m=O(N^{1/3}),
\end{equation}
in Step 5 and explains its composition-counting origin
\cite[pp.~4, 9--10]{canfield}. In particular, the sufficiency of
$N/m^3\to\infty$ is already a classical consequence. The critical leading
constant for $N=c m^3$ is explicitly posted in A238608 \cite{oeis238608}.

The polynomial part of a restricted partition function, called its first
Sylvester wave, is also classical. Dilcher and Vignat give explicit formulas
using higher-order Bernoulli polynomials \cite{dilchervignat}. Sills and
Zeilberger give an algorithmic quasipolynomial approach \cite{sillszeilberger};
O'Sullivan studies waves and their asymptotics in a different joint-growth
regime \cite{osullivan}.

The contribution here is a complete, reusable treatment of the following
more precise package: finite elementary inequalities proving necessity as
well as sufficiency; a uniform all-order theorem with all other arcs controlled;
explicit critical coefficients through order five; an expansion with errors
small enough to resolve $q^{-m}$ sectors; explicit inverse charts; and a
probabilistic boundary law with its first correction. The retrieved sources
were not found to list these particular higher coefficients or this particular
probability correction. That limited observation is not an exhaustive
bibliographic novelty certification.

\subsection{Relation to the ProveIt repository}\label{rpc:sec:proveit}

The inspected repository file
\begin{quote}\small\ttfamily
Analysis/FabiusFunction/Lean/FabiusFunction/\\
PartitionBoundedParts.lean
\end{quote}
provides a genuine formal starting point \cite{proveit}. In namespace
\texttt{Fabius}, it defines \texttt{boundedCount}, identifies restricted
partitions with multiplicity vectors through
\texttt{card\_restricted\_le\_eq}, and proves the finite-product generating
function in \texttt{hasSum\_boundedCount\_mul\_pow}. The present analytic
arguments can begin from precisely that generating function. We inspected
the source but did not run the repository's Lean build; none of the new
asymptotic or probabilistic theorems is represented as already formalized.

\subsection{Provenance and place in the collection}\label{rpc:sec:provenance}
\begin{writenote}[Added 1 October 2026, batch~74]
This report is a single manuscript, printed in full: manuscript~05 of
batch~74 of ProveIt's incoming-reports intake, delivered as the archive
\texttt{OEIS\_\allowbreak Restricted\_\allowbreak Partitions\_\allowbreak Cubic\_\allowbreak Boundary.zip} in commit
\texttt{c664fc2f0} and placed in the research-report collection, among the
OEIS-sequence asymptotics reports, in commit \texttt{b669cff87}. It names no
ProveIt commit: its repository input was a direct reading of the main-branch
file \nolinkurl{PartitionBoundedParts.lean}. At the placement commit that file
still has the namespace \texttt{Fabius} (line~22) and the declarations
\texttt{card\_restricted\_le\_eq} (line~57), \texttt{boundedCount}
(line~85) and \texttt{hasSum\_boundedCount\_mul\_pow} (line~90), as
described above. Those declarations are formal; nothing in this report is.
Placement in the collection, beside that Lean development, confers no formal
status on any statement here. The writing step changed only the labels
(every label now carries the prefix \texttt{rpc:}), added these
\textsf{[write]} notes, three labels and two references
\cite{rpc-tai,rpc-a097356}; no statement or proof was altered.

The nearest report in the collection is
\texttt{a097356-\allowbreak sqrt-\allowbreak restricted-\allowbreak partitions}
\cite{rpc-a097356}, on the same counts $p_m(N)$ at the quadratic scale
$N\asymp m^2$ (its parameter is $\alpha=N/m^2$, held in a compact set).
This report works at $N\asymp m^3$ and beyond, which is the end
$\alpha\to\infty$ of that report's research question ``Uniform limits as
the saddle moves to zero or infinity''. It does not answer that question: no
expansion here is matched uniformly to the quadratic regime, and the two
reports share no proposition. The partition reports
\texttt{a033552-\allowbreak catalan-\allowbreak partitions} (parts restricted to Catalan numbers) and
\texttt{a022629-\allowbreak distinct-\allowbreak partition-\allowbreak norms} (the products
$\prod(1+k^aq^k)$) concern different products. Sufficiency in the cubic
criterion is the classical Erd\H{o}s--Lehner estimate, as
\S\ref{rpc:sec:intro} says; the leading constant of A238608 is posted in
that entry.
\end{writenote}

\section{An elementary sharp criterion}\label{rpc:sec:elementary}

Write $R_m(N)=p_m(N)/V_m(N)$. The scale $m^3$ can be detected without
saddle points, Bernoulli numbers, or complex analysis.

\begin{lemma}[Lattice-volume bounds]\label{rpc:lem:volume}
For $m\ge1$ and $N\ge1$, put $A_m=m(m+1)/2-1$. Then
\begin{equation}\label{rpc:eq:volume}
 1\le R_m(N)\le\left(1+\frac{A_m}{N}\right)^{m-1}.
\end{equation}
\end{lemma}
\begin{proof}
The multiplicity of the part $1$ is determined by the other multiplicities.
Consequently $p_m(N)$ is the number of integer points in
\[
 \mathcal S_N=\left\{(x_2,\ldots,x_m)\in\mathbb R_{\ge0}^{m-1}:
                         \sum_{j=2}^m jx_j\le N\right\}.
\]
This simplex has volume $N^{m-1}/((m-1)!\prod_{j=2}^m j)=V_m(N)$.
Every point of $\mathcal S_N$ lies in the unit box based at its coordinatewise
floor, a feasible integer point. Conversely, every unit box based at a feasible
point lies in $\mathcal S_{N+A_m}$. The boxes have disjoint interiors and unit
volume. Taking volumes proves both bounds. For $m=1$ the statement is the
zero-dimensional identity $p_1(N)=1$.
\end{proof}

\begin{lemma}[A transposition lower bound]\label{rpc:lem:burnside}
For $m\ge3$ and $N\ge1$,
\begin{equation}\label{rpc:eq:burnside}
 R_m(N)\ge
 \prod_{r=1}^{m-1}\left(1+\frac rN\right)
 +\frac{m(m-1)^2}{4N}
 \ge 1+\frac{m(m-1)^2}{4N}.
\end{equation}
\end{lemma}
\begin{proof}
Let the symmetric group act on weak compositions
$(x_1,\ldots,x_m)$ of $N$ by permuting coordinates. Its orbits are partitions
into at most $m$ parts. Burnside's identity gives
\[
 p_m(N)=\frac1{m!}\sum_{\sigma\in S_m}|\Fix(\sigma)|.
\]
The identity fixes $\binom{N+m-1}{m-1}$ compositions, which contributes the
product in \eqref{rpc:eq:burnside} after normalization by $V_m(N)$.

A transposition requires two coordinates to be equal. Its fixed-point count
is the number of solutions of
\[
 2y+z_1+\cdots+z_{m-2}=N,\qquad y,z_i\ge0.
\]
Eliminate $z_{m-2}$ and apply the lower unit-box argument to the resulting
simplex in dimension $m-2$. The count is at least
$N^{m-2}/(2(m-2)!)$. There are $\binom m2$ transpositions. Their total
contribution, divided by $m!V_m(N)$, is $m(m-1)^2/(4N)$. All other terms
in Burnside's sum are nonnegative.
\end{proof}

\begin{theorem}[Exact threshold]\label{rpc:thm:iff}
For every positive-integer sequence $N=N(m)$,
\begin{equation}\label{rpc:eq:iff}
 p_m(N)\sim V_m(N)
 \quad\Longleftrightarrow\quad N/m^3\longrightarrow\infty.
\end{equation}
Moreover, when $m^3/N\to0$,
$R_m(N)=1+O(m^3/N)$, whereas $N/m^3\to0$ implies $R_m(N)\to\infty$.
\end{theorem}
\begin{proof}
By \eqref{rpc:eq:volume},
$0\le\log R_m(N)\le(m-1)A_m/N=O(m^3/N)$.
This proves sufficiency and its error estimate. If $R_m(N)\to1$,
\eqref{rpc:eq:burnside} forces $m(m-1)^2/N\to0$, which is equivalent to
$m^3/N\to0$. The same lower bound proves the last assertion.
\end{proof}

A useful extension of the lower bound also explains the critical exponential.
For $0\le r\le\lfloor(m-1)/2\rfloor$, retain the permutations consisting of
$r$ disjoint transpositions and $m-2r$ fixed points. The identical simplex
argument gives
\begin{equation}\label{rpc:eq:involutions}
 R_m(N)\ge
 \sum_{r=0}^{\lfloor(m-1)/2\rfloor}
 \frac{\fall m{2r}\,\fall{(m-1)}r}{4^r r!N^r}.
\end{equation}
Here $\fall xr=x(x-1)\cdots(x-r+1)$ and $\fall x0=1$.
For $N\sim c m^3$, any fixed finite truncation tends to the corresponding
truncation of $\exp(1/(4c))$. Thus
\[
 \liminf_{m\to\infty}R_m(N)\ge e^{1/(4c)}.
\]
An upper bound with the same constant needs more information; it is supplied
by the uniform expansion below. Formula \eqref{rpc:eq:involutions} is not itself
an assertion that all remaining permutation types can simply be discarded.

\section{A uniform centered expansion to every fixed order}\label{rpc:sec:centered}

\subsection{Centering and the coefficient recurrence}

Set
\begin{equation}\label{rpc:eq:center}
 S_r(m)=\sum_{j=1}^m j^r,\qquad
 M=N+\frac{S_1(m)}2=N+\frac{m(m+1)}4.
\end{equation}
The exact factorization
\begin{equation}\label{rpc:eq:factor}
 e^{Nt}F_m(e^{-t})=
 \frac{e^{Mt}}{m!t^m}\mathcal A_m(t),\qquad
 \mathcal A_m(t)=\prod_{j=1}^m\frac{jt}{2\sinh(jt/2)}
\end{equation}
is the reason for choosing this center. The remaining factor is even. Its
logarithm, in a neighborhood of zero, is
\begin{equation}\label{rpc:eq:bernoulli}
 \log\mathcal A_m(t)=\sum_{r\ge1}\gamma_r(m)t^{2r},\qquad
 \gamma_r(m)=-\frac{B_{2r}S_{2r}(m)}{2r(2r)!}.
\end{equation}
We use $B_2=1/6$, $B_4=-1/30$, and $B_6=1/42$. Formula
\eqref{rpc:eq:bernoulli} follows by differentiating
$\log(z/(2\sinh(z/2)))$ and integrating its Bernoulli expansion at zero.

Define $a_j(m)$ by $\mathcal A_m(t)=\sum_{j\ge0}a_j(m)t^{2j}$. Then
\begin{equation}\label{rpc:eq:a-rec}
 a_0=1,\qquad
 j a_j=\sum_{r=1}^j r\gamma_r a_{j-r}.
\end{equation}
The first coefficients are
\begin{align}
 a_1&=-\frac{S_2}{24},\label{rpc:eq:a1}\\
 a_2&=\frac{S_2^2}{1152}+\frac{S_4}{2880},\\
 a_3&=-\frac{S_2^3}{82944}-\frac{S_2S_4}{69120}
                         -\frac{S_6}{181440}.
\end{align}
Each $a_j$ is a rational polynomial in $m$ of degree at most $3j$.
Indeed $S_{2r}$ has degree $2r+1\le3r$, and the recurrence preserves the bound.

\begin{theorem}[Uniform centered all-order expansion]\label{rpc:thm:centered}
Fix $c_0>0$ and an integer $K\ge0$. Uniformly over integer pairs with
$N\ge c_0m^3$, as $m\to\infty$,
\begin{equation}\label{rpc:eq:centered}
 p_m(N)=\frac{M^{m-1}}{m!(m-1)!}
 \left\{\sum_{j=0}^K
  \frac{a_j(m)\fall{(m-1)}{2j}}{M^{2j}}
       +O_{K,c_0}(m^{-K-1})\right\}.
\end{equation}
In particular, the $j$th displayed term is $O_{j,c_0}(m^{-j})$.
The theorem concerns fixed truncation orders; it does not assert convergence
of the resulting infinite asymptotic series.
\end{theorem}

\subsection{The minor-arc estimate, including very large \texorpdfstring{$N$}{N}}

We give the details because a formal expansion at $q=1$ alone does not justify
\eqref{rpc:eq:centered}. Other roots of unity must be controlled uniformly while
the number of factors grows.

\begin{lemma}[Decay on a Cauchy circle]\label{rpc:lem:arcs}
Let $0<\tau\le C_0m^{-2}$, and put
\[
 \rho_m(\theta)=\frac{|F_m(e^{-\tau-\ii\theta})|}{F_m(e^{-\tau})}.
\]
For all sufficiently large $m$ there are positive constants $b,\delta,C$,
with dependence on $C_0$ allowed, such that
\begin{align}
 \rho_m(\theta)&\le
    \left(1+b(\theta/\tau)^2\right)^{-m/2},
                       &&|\theta|\le1/m,\label{rpc:eq:neararc}\\
 \rho_m(\theta)&\le(Cm\tau)^{\delta m},
                       &&1/m\le|\theta|\le\pi.\label{rpc:eq:fararc}
\end{align}
\end{lemma}
\begin{proof}
For each factor,
\[
 \frac{1-e^{-j\tau}}{|1-e^{-j\tau-\ii j\theta}|}
 =\left(1+\frac{4e^{-j\tau}\sin^2(j\theta/2)}
                  {(1-e^{-j\tau})^2}\right)^{-1/2}.
\]
If $|\theta|\le1/m$, then $|j\theta|\le1$ and $j\tau=O(1/m)$.
The inequalities $|\sin(j\theta/2)|\ge c j|\theta|$ and
$1-e^{-j\tau}\le j\tau$ prove \eqref{rpc:eq:neararc} after multiplication.

For \eqref{rpc:eq:fararc}, a fixed positive fraction of the indices $j\le m$
have $|\sin(j\theta/2)|$ bounded away from zero, uniformly in this interval.
Here is an explicit justification. For $1/m\le|\theta|\le4\pi/m$, choose
$j$ between $m/(16\pi)$ and $m/(8\pi)$; the arguments stay in a compact
subinterval of $(0,\pi)$. For $4\pi/m\le|\theta|\le\pi$, the geometric-sum
bound gives
\[
 \sum_{j=1}^m\sin^2(j\theta/2)
 \ge\frac m2-\frac1{2|\sin(\theta/2)|}\ge\frac{3m}{8}.
\]
This forces a positive fraction of the summands to exceed a fixed positive
constant. Each corresponding factor is at most $Cm\tau$; every other factor
is at most one. Taking the minimum of the constants from the two cases proves
the assertion.
\end{proof}

\subsection{Proof of the uniform theorem}
\begin{proof}[Proof of Theorem~\ref{rpc:thm:centered}]
Cauchy's formula, with $q=e^{-t}$, is
\begin{equation}\label{rpc:eq:cauchy}
 p_m(N)=\frac1{2\pi}\int_{-\pi}^{\pi}
       e^{N(\tau+\ii\theta)}F_m(e^{-\tau-\ii\theta})\dd\theta.
\end{equation}
Choose $\tau=m/M$ and the central interval
$|\theta|\le\tau m^{-2/5}$. We use three estimates.

\paragraph{Analytic Taylor remainder.}
For a sufficiently small absolute $\varepsilon>0$,
$\mathcal A_m(t)$ is analytic and uniformly bounded on
$|t|\le\varepsilon m^{-3/2}$. In fact $j|t|=O(m^{-1/2})$ there, and the
sum of the absolute logarithmic bounds is
$O(S_2(m)|t|^2)=O(1)$. Cauchy's coefficient estimate and evenness imply
\[
 \mathcal A_m(t)=\sum_{j=0}^K a_j(m)t^{2j}
        +O_K\!\left(m^{3K+3}|t|^{2K+2}\right)
\]
inside half that disk. The central interval is contained in this smaller
disk for all sufficiently large $m$, uniformly in $N\ge c_0m^3$.

\paragraph{Integrating the central remainder.}
For fixed nonnegative $r$ and $m>r+2$,
\[
 \int_{-\infty}^{\infty}|\tau+\ii\theta|^{-m+r}\dd\theta
 =\tau^{1-m+r}\sqrt\pi\,
       \frac{\Gamma((m-r-1)/2)}{\Gamma((m-r)/2)}.
\]
The gamma quotient is $O_r(m^{-1/2})$. Together with Stirling's formula,
this shows that integrating the Taylor remainder in \eqref{rpc:eq:factor},
relative to $M^{m-1}/(m!(m-1)!)$, costs at most
\[
 O_K\!\left(m^{3K+3}\tau^{2K+2}\right)
 =O_{K,c_0}(m^{-K-1}).
\]
The elementary inverse-Laplace identity
\[
 \frac1{2\pi\ii}\int_{\tau-\ii\infty}^{\tau+\ii\infty}
       e^{Mt}t^{-d}\dd t=\frac{M^{d-1}}{(d-1)!},\qquad d\ge2,
\]
follows by closing the contour to the left and taking the residue at zero.
Apply it with $d=m-2j$. Extending each central polynomial integral to this
line changes it by an exponentially small quantity: the gamma kernel outside
$|\theta|\le\tau m^{-2/5}$ has a tail bounded by
$\exp(-b_1m^{1/5})$ times its natural integral size. This gives exactly the
falling factorial in \eqref{rpc:eq:centered}.

\paragraph{Discarding the rest of the original circle.}
On $\tau m^{-2/5}\le|\theta|\le1/m$, \eqref{rpc:eq:neararc} and the change
of variable $v=\theta/\tau$ give a relative contribution
$O(\exp(-b_2m^{1/5}))$. On the remaining arcs, the relative contribution is
bounded by
\begin{equation}\label{rpc:eq:far-prefactor}
 C\frac{\sqrt m}{\tau}(Cm\tau)^{\delta m}.
\end{equation}
The factor $\sqrt m/\tau$ comes from normalizing the real-axis integrand by
its gamma-kernel integral. It cannot be ignored when $N$ is extremely large.
For $\delta m>1$, the right side is increasing in $\tau$. Substituting
$\tau\le C_0/m^2$ bounds it by
$m^{5/2}(C_1/m)^{\delta m}$, which is smaller than every fixed inverse power
of $m$. Thus the argument is uniform even for factorial or faster growth of
$N$, rather than merely for $N$ bounded by a power of $m$.

Combining the three estimates proves \eqref{rpc:eq:centered}.
\end{proof}

\begin{remark}[What this theorem does not assert]
For fixed $m$, the polynomial part of $p_m(N)$ is its first Sylvester wave.
The coefficient algebra above is consistent with that classical polynomial.
The theorem gives uniform finite-order asymptotics of the entire count, not
just of a polynomial formally extracted at one singularity. It does not
identify each exponentially small root-of-unity wave, provide a Stokes
analysis, or specify an optimal truncation order growing with $m$.
\end{remark}

\section{Critical constants and explicit OEIS expansions}\label{rpc:sec:critical}

\begin{corollary}[Critical crossover]\label{rpc:cor:crossover}
If $N/m^3\to c\in(0,\infty)$, then
\begin{equation}\label{rpc:eq:crossover}
 \frac{p_m(N)}{V_m(N)}\longrightarrow e^{1/(4c)}.
\end{equation}
\end{corollary}
\begin{proof}
The $K=0$ case of Theorem~\ref{rpc:thm:centered} is
$p_m(N)=M^{m-1}/(m!(m-1)!)(1+O(1/m))$.
Furthermore,
\[
 (m-1)\log(M/N)
 =(m-1)\log\left(1+\frac{m(m+1)}{4N}\right)\longrightarrow\frac1{4c}.
\]
\end{proof}

\subsection{A general all-order logarithmic expansion}

Write $c=N/m^3$; this notation is permitted even when $c$ varies with $m$.
For every fixed $K$, uniformly for $c\ge c_0>0$,
\begin{align}\label{rpc:eq:critical-log}
 \log p_m(c m^3)
 ={}&(m-3)\log m+2m+(m-1)\log c-\log(2\pi)+\frac1{4c}\notag\\
 &+\sum_{r=1}^K\frac{L_r(c)}{m^r}+O_{K,c_0}(m^{-K-1}).
\end{align}
Here and below the left side is used only when $c m^3$ is an integer.
The $L_r(c)$ are rational polynomials in $1/c$. Equivalently,
\begin{equation}\label{rpc:eq:critical-relative}
 p_m(c m^3)=\frac{e^{2m+1/(4c)}c^{m-1}m^{m-3}}{2\pi}
 \left(\sum_{r=0}^K\frac{b_r(c)}{m^r}+O_{K,c_0}(m^{-K-1})\right),
\end{equation}
where
\begin{equation}\label{rpc:eq:b-rec}
 b_0=1,\qquad k b_k=\sum_{r=1}^k rL_r b_{k-r}.
\end{equation}

To obtain \eqref{rpc:eq:critical-log}, take the logarithm of
\eqref{rpc:eq:centered}, substitute $M=c m^3(1+(m^{-1}+m^{-2})/(4c))$, and
use the classical Stirling expansion \cite{dlmf}:
\begin{equation}\label{rpc:eq:stirling}
 \log\Gamma(m+1)+\log\Gamma(m)
 \sim 2m\log m-2m+\log(2\pi)
       +2\sum_{r\ge1}\frac{B_{2r}}{2r(2r-1)m^{2r-1}}.
\end{equation}
Every operation is justified to a fixed order by Theorem~\ref{rpc:thm:centered}.
Thus the coefficient algorithm is backed by an asymptotic remainder theorem,
not only by formal manipulation.

The first five logarithmic coefficients are
\begin{align}
 L_1(c)&=-\frac{48c^2+13}{288c^2},\label{rpc:eq:L1}\\
 L_2(c)&=-\frac{144c^2+6c-7}{576c^3},\label{rpc:eq:L2}\\
 L_3(c)&=\frac{11520c^4+122400c^2+14400c-8081}{2073600c^4},\\
 L_4(c)&=\frac{21600c^3-50400c^2-7855c+2861}{2073600c^5},\\
 L_5(c)&=-\frac{1}{21946982400c^6}\bigl(
 34836480c^6+304819200c^4+304819200c^3\notag\\[-2pt]
 &\hspace{40mm}{}-235282320c^2-42928704c+11425901\bigr).
\end{align}
As a short independent check of the first coefficient, the centered shift
contributes $-1/(32c^2m)$, the $a_1$ term contributes
$-1/(72c^2m)$, and Stirling contributes $-1/(6m)$.
Their sum is \eqref{rpc:eq:L1}.

\subsection{Five corrections for A238608}

At $c=1$, the result becomes
\begin{align}\label{rpc:eq:A238608}
 p_m(m^3)=\frac{e^{2m+1/4}m^{m-3}}{2\pi}\Bigg[
 &1-\frac{61}{288m}-\frac{37463}{165888m^2}
 +\frac{425075267}{3583180800m^3}\notag\\
 &-\frac{21833043287}{4127824281600m^4}
 -\frac{2048618007321713}{58251856261939200m^5}
 +O(m^{-6})\Bigg].
\end{align}
The leading factor is the one already posted in A238608
\cite{oeis238608}. The additional terms and their remainder are the refinement.
The logarithmic form is often more stable numerically:
\begin{align}\label{rpc:eq:c1log}
 \log\!\left(\frac{2\pi p_m(m^3)}{e^{2m+1/4}m^{m-3}}\right)
 ={}&-\frac{61}{288m}-\frac{143}{576m^2}
 +\frac{140239}{2073600m^3}\notag\\
 &-\frac{16897}{1036800m^4}
 -\frac{377689757}{21946982400m^5}+O(m^{-6}).
\end{align}
Substituting $c=2,3,4,5$ in the general formulas immediately gives the same
order of refinement for A258302--A258305, without separate saddle analyses.

\subsection{Higher power columns and an important rounding issue}

For any fixed integer $k\ge4$, substitute $c=m^{k-3}$ in
\eqref{rpc:eq:critical-log}. One obtains the complete inverse-power expansion of
\begin{equation}\label{rpc:eq:higherpower}
 p_m(m^k)\sim\frac{e^{2m}m^{(k-2)m-k}}{2\pi},
\end{equation}
the leading expression posted in the A238016 family \cite{oeis238016}.
For example, the quartic column A238609 satisfies
\begin{align}\label{rpc:eq:quartic}
 \log p_m(m^4)={}&(2m-4)\log m+2m-\log(2\pi)
       +\frac1{12m}-\frac{139}{480m^3}\notag\\
       &-\frac1{96m^4}+\frac{1403}{20160m^5}+O(m^{-6}).
\end{align}
In particular, its first correction has the opposite sign from the cubic
column's first correction.

If the intended count is $p_m(\lfloor c m^3\rfloor)$ for a fixed nonintegral
$c$, use $c_m=\lfloor c m^3\rfloor/m^3$ in \eqref{rpc:eq:critical-log}.
Replacing $c_m$ indiscriminately by $c$ changes the logarithm by $O(m^{-2})$,
which is large enough to corrupt coefficients beginning at the second
correction. This is a discrete rounding effect, not a failure of the theorem.

\section{Dilute expansions and genuine exponential sectors}\label{rpc:sec:dilute}

The centered theorem is uniform in very large $N$, but its remainder
$O(m^{-K-1})$ is not small enough to detect a term such as $2^{-m}$.
A separate expansion in $1/N$, with a remainder measured in $m^3/N$, is needed.

\subsection{The uncentered coefficient system}

Define
\begin{equation}\label{rpc:eq:H}
 \mathcal H_m(t)=\prod_{j=1}^m\frac{jt}{1-e^{-jt}}
 =\exp\left(\frac{S_1t}{2}+\sum_{r\ge1}\gamma_rt^{2r}\right)
 =\sum_{j\ge0}h_j(m)t^j.
\end{equation}
Put
\begin{equation}\label{rpc:eq:beta}
 \beta_j(m)=\fall{(m-1)}j h_j(m),\qquad \beta_0=1.
\end{equation}
If $\lambda_1=S_1/2$, $\lambda_{2r}=\gamma_r$, and
$\lambda_{2r+1}=0$ for $r\ge1$, then
\begin{equation}\label{rpc:eq:h-rec}
 jh_j=\sum_{r=1}^j r\lambda_r h_{j-r}.
\end{equation}
Thus $h_j$ and $\beta_j$ have rational coefficients and degrees at most $2j$
and $3j$, respectively.

\begin{theorem}[Uniform dilute expansion]\label{rpc:thm:dilute}
There is an absolute $\eta_0>0$ such that, for every fixed $K\ge0$, uniformly
for $0<\eta=m^3/N\le\eta_0$ and all sufficiently large $m$,
\begin{equation}\label{rpc:eq:dilute}
 \frac{p_m(N)}{V_m(N)}
   =\sum_{j=0}^K\frac{\beta_j(m)}{N^j}
      +O_K\!\left((m^3/N)^{K+1}\right).
\end{equation}
In particular, this holds along every sequence for which $m^3/N\to0$.
\end{theorem}
\begin{proof}
Use \eqref{rpc:eq:cauchy} with $\tau=m/N=\eta/m^2$. The exact integrand now is
$e^{Nt}\mathcal H_m(t)/(m!t^m)$. For a sufficiently small fixed $r>0$,
$\mathcal H_m$ is uniformly bounded on $|t|\le r/m^2$: the linear logarithmic
term is bounded, and all other terms have total absolute size $O(m^3|t|^2)$.
Cauchy's estimate gives
\[
 \mathcal H_m(t)=\sum_{j=0}^K h_jt^j
                   +O_K((m^2|t|)^{K+1})
\]
on the disk of half this radius. Choose $\eta_0$ sufficiently small that
$\tau$ lies well inside the disk and integrate over
$|\theta|\le r/(4m^2)$. The same gamma-kernel calculation as before bounds
the relative remainder by
$(m^2\tau)^{K+1}=\eta^{K+1}$. Each polynomial term gives
$h_j\fall{(m-1)}j/N^j$ by inverse Laplace integration.

For completeness, the omitted arcs are small on this stronger scale, not
merely small in $1/m$. On $r/(4m^2)\le|\theta|\le1/m$,
Lemma~\ref{rpc:lem:arcs}, after $v=\theta/\tau$, bounds the normalized tails
by a polynomial factor in $m$ times $(C\eta)^{m-1}$. For sufficiently small
$\eta_0$ and $m$ sufficiently large relative to $K$, this is
$O_K(\eta^{K+1})$. On $1/m\le|\theta|\le\pi$, the analogous bound is
\[
 C\frac{m^{5/2}}{\eta}\left(\frac{C\eta}{m}\right)^{\delta m}.
\]
After division by $\eta^{K+1}$, the positive remaining exponent of $\eta$
allows substitution of $\eta_0$, leaving a quantity tending rapidly to zero.
These bounds also justify extending the polynomial integrals to the entire
vertical line. This proves \eqref{rpc:eq:dilute}.
\end{proof}

Taking the formal logarithm defines rational polynomials $D_j(m)$ through
\begin{equation}\label{rpc:eq:D-rec}
 D_j=\beta_j-\frac1j\sum_{r=1}^{j-1}rD_r\beta_{j-r}.
\end{equation}
Theorem~\ref{rpc:thm:dilute} justifies the corresponding logarithmic expansion
with the same order of remainder. Its first terms are
\begin{align}
 D_1(m)&=\frac{m(m^2-1)}4,\label{rpc:eq:D1}\\
 D_2(m)&=-\frac{m(m^2-1)(13m^2+3m-4)}{288}.\label{rpc:eq:D2}
\end{align}
For example $h_2=S_1^2/8-S_2/24$, so
$D_2=(m-1)(m-2)h_2-D_1^2/2$, an immediate exact check.

\subsection{Exponential and superexponential specializations}

For a fixed integer $q\ge2$, put $a_q(m)=p_m(q^m)$ and $b=\log q$.
The full fixed-order expansion is
\begin{align}\label{rpc:eq:exp-transseries}
 \log a_q(m)={}&b m(m-1)-\log\Gamma(m+1)-\log\Gamma(m)\notag\\
 &+\sum_{j=1}^K D_j(m)e^{-jbm}
      +O_K\!\left(m^{3K+3}e^{-(K+1)bm}\right).
\end{align}
This is an expansion in genuine exponentially separated sectors, rather than
only a power series in $1/m$ renamed a transseries. The exact gamma background
is retained so that its algebraic truncation error does not obscure the
exponential corrections. In particular,
\begin{align}\label{rpc:eq:exp-two}
 \log a_q(m)={}&b m(m-1)-\log\Gamma(m+1)-\log\Gamma(m)\notag\\
 &+\frac{m(m^2-1)}4q^{-m}
 -\frac{m(m^2-1)(13m^2+3m-4)}{288}q^{-2m}
 +O(m^9q^{-3m}).
\end{align}
The leading term corresponds to A238010 and its columns, including A237998
when $q=2$ \cite{oeis238010}.

For A238000, put $N=m^m$ in Theorem~\ref{rpc:thm:dilute}. The sectors are
$m^{-jm}=\exp(-jm\log m)$, with the same polynomial coefficients $D_j(m)$.
For A258670, put $N=(2m)!$; its sectors are $((2m)!)^{-j}$
\cite{oeis258670}. In particular,
\[
 p_m((2m)!)=\frac{((2m)!)^{m-1}}{m!(m-1)!}
 \left(1+\frac{m(m^2-1)}{4(2m)!}
       +O\left(\frac{m^6}{((2m)!)^2}\right)\right).
\]
There is no need to solve a new saddle equation for each such sequence.
The full collection follows from one uniform dilute theorem.

\section{Inverting the asymptotic expansions}\label{rpc:sec:inverse}

\subsection{What inversion means for an integer sequence}

A sequence has no canonical smooth interpolation. We therefore use an
\emph{asymptotic inverse chart}: a smooth expression $I(y)$ that, at the exact
sample values $y=a(m)$, recovers $m$ with a specified asymptotic error.
The same expression can approximate an integer threshold inverse
$Q(y)=\min\{m:a(m)\ge y\}$, but generally only up to an $O(1)$ rounding
ambiguity. An exponentially accurate sample-point chart does not magically
supply an exponentially accurate integer-valued function for all real $y$.

Each fixed truncation of a forward logarithmic profile is smooth for large
real $x$ and can be inverted there. Raising the truncation order produces
consistent inverse coefficients. At sample values the forward remainder,
divided by the derivative of the leading profile, controls the inverse error.
This is the precise sense of the all-order inversions used below.

\begin{writenote}[Added 1 October 2026, batch~74: an instance of repository results]
The distinction drawn here between a sample-point chart and the integer
threshold $Q(y)$, with its $O(1)$ rounding ambiguity, is the content of the
staircase theorem \texttt{p0:thm:staircase} of the transseries-and-inversion
volume \cite{rpc-tai}: $Q=\lceil\nu\rceil$ for any increasing
interpolation $\nu$, and a chart with error at most $\eta$ determines $Q$
exactly when the separation condition of its part~(2) holds. The
inversions of this section are instances of that volume's inversion
apparatus, and no novelty is claimed for the method.
\end{writenote}

\subsection{The cubic family: Lambert \texorpdfstring{$W$}{W} and two explicit corrections}

Fix $c>0$, and consider $a_c(m)=p_m(c m^3)$ on admissible integer arguments.
Write
\[
 b_c=2+\log c,\qquad D_c=\frac1{4c}-\log(2\pi c).
\]
Equation \eqref{rpc:eq:critical-log} has the form
\begin{equation}\label{rpc:eq:cubicprofile}
 \log a_c(m)=m\log m+b_cm-3\log m+D_c
               +\frac{L_1(c)}m+\frac{L_2(c)}{m^2}+\cdots.
\end{equation}
For large $y$, set $L=\log y$ and
\begin{equation}\label{rpc:eq:lambertbase}
 w=W(c e^2L),\qquad x_0=\frac Lw.
\end{equation}
Here $W$ is the positive real solution of $we^w=z$. Then
$\log x_0+b_c=w$, and consequently
$x_0\log x_0+b_cx_0=L$ exactly. Define
\begin{equation}\label{rpc:eq:deltas}
 d_0=\frac{3\log x_0-D_c}{w+1},\qquad
 d_1=\frac{3d_0-d_0^2/2-L_1(c)}{w+1}.
\end{equation}

\begin{theorem}[Cubic inverse chart]\label{rpc:thm:cubicinverse}
At $y=a_c(m)$, with $c$ fixed,
\begin{equation}\label{rpc:eq:inverse-cubic}
 m=x_0+d_0+\frac{d_1}{x_0}+O_c(x_0^{-2}).
\end{equation}
There is an all-order expansion
$x_0+\sum_{r\ge0}d_r(w;c)x_0^{-r}$, in the fixed-order asymptotic sense.
Its coefficients are obtained recursively by Taylor substitution in
\eqref{rpc:eq:cubicprofile}; they are rational functions of $w$ with constants
depending on $c$.
\end{theorem}
\begin{proof}
Let $f_0(x)=x\log x+b_cx$. At $x_0$, one has
$f_0'=w+1$ and $f_0''=1/x_0$. Substitute
$x=x_0+d_0+d_1/x_0$ in \eqref{rpc:eq:cubicprofile}.
The constant residual is
$(w+1)d_0-3\log x_0+D_c$, and the coefficient of $x_0^{-1}$ is
$(w+1)d_1+d_0^2/2-3d_0+L_1(c)$. Both vanish by \eqref{rpc:eq:deltas}.
The remaining residual is $O_c(x_0^{-2})$, whereas the derivative is
asymptotic to $w+1$. The mean-value theorem proves the stated, slightly
weaker, $O_c(x_0^{-2})$ inverse error.

At any further fixed order, the new unknown coefficient appears multiplied
by $w+1$; everything else is already known. Division by $w+1$ completes
an induction. Since $\log x_0=w-b_c$, the recursion stays within rational
functions of $w$. The forward expansion can be taken to arbitrarily high
fixed order by Theorem~\ref{rpc:thm:centered}.
\end{proof}

For A238608 specifically,
\[
 d_0=3-\frac{37/4-\log(2\pi)}{w+1},\qquad
 d_1=\frac{3d_0-d_0^2/2+61/288}{w+1},\qquad w=W(e^2\log y).
\]
These corrections are important even for moderately large indices: the raw
Lambert term can miss the index by nearly two while the corrected expression
is already accurate to about $10^{-4}$ at $m=60$.

\begin{writenote}[Added 1 October 2026, batch~74: an instance of repository results]
The core $x\log x+b_cx=L$, solved by the principal branch as in
\eqref{rpc:eq:lambertbase}, is the Lambert core of the inverse gamma function
in \texttt{p6:thm:gamma} of the transseries-and-inversion volume
\cite{rpc-tai} (there $T\log T-T=R$ with $T=R/W_0(R/e)$), shifted by the
constant $b_c$; the corrections $d_0,d_1,\dots$ are its recursive
perturbed inversion \texttt{p0:thm:perturbed-inversion}. The exponential
columns below use the same perturbed-inversion step with the small parameter
$e^{-bX}$. Theorem~\ref{rpc:thm:cubicinverse} and \eqref{rpc:eq:expinverse} are
instances of that apparatus.
\end{writenote}

\subsection{Exponential columns: algebraic and invisible corrections}

For $q\ge2$ fixed, define the exact smooth background
\begin{equation}\label{rpc:eq:F0}
 G_q(x)=b x(x-1)-\log\Gamma(x+1)-\log\Gamma(x),\qquad b=\log q.
\end{equation}
Stirling's expansion yields
\[
 G_q(x)=b x^2-2x\log x+(2-b)x-\log(2\pi)
                -\frac1{6x}+\frac1{180x^3}-\cdots.
\]
Put
\begin{equation}\label{rpc:eq:sde}
 s=\sqrt{L/b},\qquad
 d=\frac{\log s}{b}+\frac{b-2}{2b},\qquad
 e=\frac{d^2}{2}+\frac db+\frac{\log(2\pi)}{2b}.
\end{equation}
Taylor substitution gives the algebraic inverse chart
\begin{equation}\label{rpc:eq:expalginverse}
 m=s+d+\frac es+O_q\!\left(\frac{(\log s)^3}{s^2}\right),
                         \qquad y=a_q(m).
\end{equation}
For example, the coefficient of $s$ cancels because
$2bd=2\log s+b-2$, and the constant residual cancels because
$2be=bd^2+2d+\log(2\pi)$.

To expose an exponentially small inverse correction, let $X$ be the exact
large solution of $G_q(X)=\log y$. This background inverse is well defined
for large $y$ since
\[
 G_q'(X)=b(2X-1)-\psi(X+1)-\psi(X)\sim2bX>0.
\]
Equation \eqref{rpc:eq:exp-transseries} and one implicit Taylor step give
\begin{equation}\label{rpc:eq:expinverse}
 m=X-\frac{D_1(X)e^{-bX}}{G_q'(X)}
             +O_q(X^5e^{-2bX}),\qquad y=a_q(m).
\end{equation}
The sign is negative: the positive count correction raises $\log a_q$ above
the gamma background, so the corrected index at fixed $y$ is smaller.
The $O(m^6e^{-2bm})$ forward error divided by $G_q'(m)\asymp m$ gives
the displayed safe inverse error; the Taylor cross terms are smaller.

All further exponential sectors can be computed recursively. For example,
write $m=X+u_1(X)e^{-bX}+u_2(X)e^{-2bX}+\cdots$. Then
\begin{align}\label{rpc:eq:inverse-sector-rec}
 u_1&=-D_1/G_q',\\
 u_2&=-\frac{D_2+(D_1'-bD_1)u_1+G_q''u_1^2/2}{G_q'}.
\end{align}
Theorem~\ref{rpc:thm:dilute} supplies the required forward remainder at every
fixed sector order. Retaining only finitely many terms of Stirling in
$G_q$ would not justify \eqref{rpc:eq:expinverse}: that algebraic error would
swamp the quantity being resolved.

\subsection{The diagonal sequence A238000}

The same procedure gives a particularly simple leading inverse for A238000.
Here
\[
 \log p_m(m^m)=(m^2-3m)\log m+2m-\log(2\pi)+O(1/m),
\]
with exponentially smaller corrections in $m^{-m}$ supplied by
Section~\ref{rpc:sec:dilute}. If
\[
 x_0=\sqrt{\frac{2L}{W(2L)}},
\]
then $x_0^2\log x_0=L$, and a first correction is
\begin{equation}\label{rpc:eq:diaginverse}
 m=x_0+\frac{3\log x_0-2}{2\log x_0+1}+O(x_0^{-1})
                         \quad\text{at }y=p_m(m^m).
\end{equation}
This follows by dividing the leading residual $-3x_0\log x_0+2x_0$
by the derivative $x_0(2\log x_0+1)$. It provides another inverse chart from
the same partition theorem, rather than an unrelated inversion example.

\section{A Poisson boundary law in the conjugate partition model}
\label{rpc:sec:poisson}

\subsection{Which repeated parts are being counted}

Choose a partition of $N$ into \emph{at most $m$ parts} uniformly, and pad it
with zeros to obtain
\[
 0\le x_1\le\cdots\le x_m,\qquad x_1+\cdots+x_m=N.
\]
This is the conjugate of the original bounded-part-size model.
The distinction matters: partitions whose part sizes are bounded by $m$
typically have many repeated small parts, and the theorem below is not a
claim about those multiplicities.

Set $d_1=x_1$ and $d_i=x_i-x_{i-1}$ for $i\ge2$. Then
\[
 \sum_{i=1}^m(m-i+1)d_i=N.
\]
Let
\[
 Z_m=\#\{i:2\le i\le m,\ d_i=0\},
\]
the number of equal adjacent pairs in the padded partition. A gap equal to
zero corresponds to a zero multiplicity for one of the weights $1,\ldots,m-1$
in the generating product. Its probability generating function is exactly
\begin{equation}\label{rpc:eq:pgf}
 \E[u^{Z_m}]=
 \frac{[q^N]F_m(q)\prod_{j=1}^{m-1}\bigl(u+(1-u)q^j\bigr)}{p_m(N)}.
\end{equation}
Indeed for each marked weight the sum over a gap is
$u+q^j+q^{2j}+\cdots=(u+(1-u)q^j)/(1-q^j)$.

\subsection{A deletion estimate and the limit distribution}

\begin{lemma}[Deleting a fixed number of factors]\label{rpc:lem:delete}
Fix $k$. Uniformly over $k$-element sets $S\subseteq\{1,\ldots,m\}$ and
$N/m^3$ in a compact subinterval of $(0,\infty)$,
\begin{equation}\label{rpc:eq:delete}
 \frac{[q^N]F_m(q)\prod_{j\in S}(1-q^j)}{p_m(N)}
 =\left(\prod_{j\in S}j\right)
       \frac{\fall{(m-1)}k}{M^k}\left(1+O_k(m^{-1})\right).
\end{equation}
\end{lemma}
\begin{proof}
On the central interval of Section~\ref{rpc:sec:centered},
\[
 \prod_{j\in S}(1-e^{-jt})
 =\left(\prod_{j\in S}j\right)t^k(1+O_k(m|t|))
 =\left(\prod_{j\in S}j\right)t^k(1+O_k(m^{-1})).
\]
Inserting this factor in the contour proof changes $t^{-m}$ to $t^{-m+k}$,
so the inverse-Laplace integral contributes $\fall{(m-1)}k/M^k$.
Both $\mathcal A_m(t)$ and the denominator count have relative error
$O(1/m)$ at this precision. On the other arcs the new product has modulus
at most $2^k$. Since $k$ is fixed and $\tau\asymp m^{-2}$, normalizing by
the smallest possible main deleted-factor contribution loses only a fixed
power of $m$, harmless compared with the arc decay. All bounds are uniform
in the choice of $S$.
\end{proof}

\begin{theorem}[Poisson repeated-part boundary]\label{rpc:thm:poisson}
If $N/m^3\to c\in(0,\infty)$, then
\begin{equation}\label{rpc:eq:poisson}
 Z_m\ \xrightarrow{\ d\ }\ \Pois\!\left(\frac1{2c}\right).
\end{equation}
With probability tending to one, the partition has exactly $m$ positive
parts and no part occurs three or more times. Thus $Z_m$ also agrees, with
probability tending to one, with the number of distinct positive part sizes
that occur twice.
\end{theorem}
\begin{proof}
For every fixed $k$, the factorial-moment identity from \eqref{rpc:eq:pgf} and
Lemma~\ref{rpc:lem:delete} give
\[
 \E[\fall{Z_m}k]
 =k!\,\frac{\fall{(m-1)}k}{M^k}
   \sum_{\substack{S\subseteq\{1,\ldots,m-1\}\\|S|=k}}
       \prod_{j\in S}j\,\left(1+O_k(m^{-1})\right).
\]
For fixed $k$, the elementary symmetric sum here equals
\[
 \frac1{k!}\left(\sum_{j=1}^{m-1}j\right)^k(1+O_k(m^{-1})).
\]
To see this, expand the power and remove tuples with a repeated index; their
relative total is bounded by a constant times
$\sum j^2/(\sum j)^2=O(1/m)$. Therefore
\[
 \E[\fall{Z_m}k]\longrightarrow(1/(2c))^k.
\]
These are the factorial moments of the asserted Poisson law. For clarity,
the first moment gives tightness; the next higher moments give uniform
integrability of every fixed power along subsequences. Any subsequential
limit therefore has the Poisson moments. The Poisson distribution is
moment-determinate, as its moment generating function is finite in a
neighborhood of zero. This proves \eqref{rpc:eq:poisson}.

The event $d_1=0$ is a single deletion of weight $m$, and hence has probability
$m(m-1)/M\,(1+O(1/m))=O(1/m)$. A triple repetition requires two consecutive
gaps to vanish. By \eqref{rpc:eq:delete}, summing these two-deletion probabilities
over at most $m$ adjacent pairs costs
\[
 O\left(\frac{m^2}{M^2}\sum_{j=1}^{m-1}j(j+1)\right)=O(m^{-1}).
\]
Outside these two exceptional events there are $m$ positive parts and only
double repetitions, proving the last assertion.
\end{proof}

\subsection{An explicit correction to the Poisson generating function}

The limit can be refined without assuming independence of the gaps.

\begin{theorem}[First finite-size probability correction]\label{rpc:thm:pgf}
Let $c=N/m^3$ range over a compact subinterval of $(0,\infty)$. Uniformly
for real $0\le u\le1$,
\begin{equation}\label{rpc:eq:pgf-correction}
 \E[u^{Z_m}]=e^{(u-1)/(2c)}
 \left[1-\frac1m\left(\frac{u-1}{c}
             +\frac{7u(u-1)}{24c^2}\right)+O_c(m^{-2})\right].
\end{equation}
The constant in the error can be chosen uniformly on the indicated compact
$c$ interval. In particular,
\begin{equation}\label{rpc:eq:no-repeats}
 \Prob(Z_m=0)=e^{-1/(2c)}\left(1+\frac1{cm}+O_c(m^{-2})\right).
\end{equation}
\end{theorem}
\begin{proof}
Put $v=u-1$, $S_r'=\sum_{j=1}^{m-1}j^r$, and $M_v=M+vS_1'$.
Near zero,
\[
 \log\prod_{j=1}^{m-1}\bigl(u+(1-u)e^{-jt}\bigr)
 =vS_1't-\frac{v+v^2}{2}S_2't^2+O(m^4|t|^3),
\]
uniformly for $0\le u\le1$. After combining this with
\eqref{rpc:eq:bernoulli}, define
\[
 C_v=\frac{S_2}{24}+\frac{v+v^2}{2}S_2'.
\]
The same centered contour proof, now at $\tau=m/M_v$, yields for the numerator
of \eqref{rpc:eq:pgf}
\[
 \frac{M_v^{m-1}}{m!(m-1)!}
 \left(1-C_v\frac{(m-1)(m-2)}{M_v^2}+O(m^{-2})\right).
\]
The error follows from $m^4|t|^3+m^6|t|^4=O(m^{-2})$ on the central
scale. The old arc estimates remain valid: for $0\le u\le1$ and
$|q|\le1$, every multiplier $u+(1-u)q^j$ has modulus at most one.
Also $M_v=M+O(m^2)>0$ and the change of real-axis normalization costs only
a bounded factor.

Divide by the same formula at $v=0$ and take a logarithm. With $N=cm^3$,
\begin{align*}
 (m-1)\log(M_v/M)
 &=\frac v{2c}-\frac1m\left(\frac vc+\frac{v+v^2}{8c^2}\right)
         +O(m^{-2}),\\
 -C_v\frac{(m-1)(m-2)}{M_v^2}
 +C_0\frac{(m-1)(m-2)}{M^2}
 &=-\frac{v+v^2}{6c^2m}+O(m^{-2}).
\end{align*}
Since $1/8+1/6=7/24$, exponentiation proves \eqref{rpc:eq:pgf-correction}.
Setting $u=0$ proves \eqref{rpc:eq:no-repeats}.
\end{proof}

There is a useful independent exact identity at $u=0$. Subtract the staircase
$(0,1,\ldots,m-1)$ from a strictly increasing nonnegative partition. The result
is an arbitrary weakly increasing nonnegative partition of
$N-T$, where $T=m(m-1)/2$. Thus
\begin{equation}\label{rpc:eq:staircase}
 \Prob(Z_m=0)=\frac{p_m(N-T)}{p_m(N)}.
\end{equation}
Inserting the centered expansion into this ratio checks both the limit
$e^{-1/(2c)}$ and the first correction $1/(cm)$.
Higher pgf corrections can be generated by retaining further terms in
$\log(u+(1-u)e^{-z})$ and using the same coefficient recursions. We do not
claim a total-variation rate or a uniform complex-$u$ expansion from the
real interval estimate alone.

\section{Exact computation and reproducible diagnostics}\label{rpc:sec:checks}

\subsection{Independent coefficient extraction}

The package's \texttt{verify.py} computes integer coefficients by the
standard product update
\[
 P_0(0)=1,\quad P_0(n)=0\ (n>0),\qquad
 P_j(n)=P_{j-1}(n)+P_j(n-j),
\]
where negative arguments contribute zero. A second implementation uses the
logarithmic derivative of \eqref{rpc:eq:gf}:
\begin{equation}\label{rpc:eq:independentrec}
 n p_m(n)=\sum_{r=1}^n
  \left(\sum_{\substack{d\mid r\\d\le m}}d\right)p_m(n-r).
\end{equation}
These two methods agree for $1\le m\le12$ and $0\le n\le200$.
The finite inequalities in Section~\ref{rpc:sec:elementary} are tested after
clearing denominators, so their verification uses exact integer arithmetic.
All displayed initial A238608 values through $m=16$ agree with the inspected
OEIS entry \cite{oeis238608}.

The complete run used 90-decimal-digit floating-point arithmetic only after
the exact counts were produced. It performed 9,230 exact assertions and
recorded 258 count samples, with $m\le60$ and coefficient arrays extending
to $N=320000$. The sample families include $c m^3$ for $1\le c\le5$,
$m^4$, and $2^m$. These are finite tests; the all-order and uniform statements
rest on the contour estimates, not on extrapolating the tables.

\subsection{Accuracy of the cubic expansion}

Let $A_K(m)$ be the right side of \eqref{rpc:eq:A238608}, truncated after
$m^{-K}$. The following entries are signed relative errors
$A_K(m)/p_m(m^3)-1$ from the included exact computation.
\begin{center}
\small
\begin{tabular}{@{}r r r r r@{}}
\toprule
$m$ & $K=0$ & $K=1$ & $K=3$ & $K=5$\\
\midrule
10 & $2.38780\cdot10^{-2}$ & $2.19166\cdot10^{-3}$ & $8.72320\cdot10^{-7}$ & $-2.93140\cdot10^{-8}$\\
20 & $1.12656\cdot10^{-2}$ & $5.55991\cdot10^{-4}$ & $4.40927\cdot10^{-8}$ & $-4.51323\cdot10^{-10}$\\
40 & $5.46413\cdot10^{-3}$ & $1.40056\cdot10^{-4}$ & $2.41567\cdot10^{-9}$ & $-7.04814\cdot10^{-12}$\\
60 & $3.60523\cdot10^{-3}$ & $6.24068\cdot10^{-5}$ & $4.54363\cdot10^{-10}$ & $-6.18667\cdot10^{-13}$\\
\bottomrule
\end{tabular}
\end{center}
The full CSV includes $K=2$, the normalized exact count, and the other values
of $c$. A small error at one index alone would not prove an expansion;
these data are useful as stringent checks on coefficient signs, normalization,
and omitted shifts.

\subsection{Probability and inverse diagnostics}

For $c=1$, the limiting mean is $1/2$, the limiting zero probability is
$e^{-1/2}\approx0.60653066$, and the limiting pgf at $u=1/2$ is
$e^{-1/4}\approx0.77880078$. Exact finite values are
\begin{center}
\begin{tabular}{@{}r r r r@{}}
\toprule
$m$ & $\E Z_m$ & $\Prob(Z_m=0)$ & $\E[2^{-Z_m}]$\\
\midrule
10 & 0.3846890512 & 0.6682085339 & 0.8213708910\\
20 & 0.4390680097 & 0.6371577383 & 0.8006057040\\
40 & 0.4686467899 & 0.6217743594 & 0.7898304769\\
60 & 0.4788932680 & 0.6166760199 & 0.7861818675\\
\bottomrule
\end{tabular}
\end{center}
The first-corrected pgf at $u=1/2$ has relative error about
$7.04\cdot10^{-5}$ at $m=60$, consistent with its $O(m^{-2})$ remainder.
The exact pgf there is calculated using the coefficients of
$\prod_{j=1}^{m-1}(1+q^j)$, divided by $2^{m-1}$, in \eqref{rpc:eq:pgf}.

At $y=p_{60}(60^3)$, the uncorrected Lambert chart has index error
$-1.96028$, while \eqref{rpc:eq:inverse-cubic} has error
$1.66805\cdot10^{-4}$. At $y=p_{18}(2^{18})$, the exact gamma-background
inverse $X$ has index error $2.99966\cdot10^{-4}$; adding the exponential
correction in \eqref{rpc:eq:expinverse} reduces it to
$-1.72742\cdot10^{-8}$. These diagnostics illustrate why algebraic and
exponential inverse corrections should be treated separately.

\subsection{Reproduction and package contents}

The archive contains the article source and PDF, plus:
\begin{description}[style=nextline,leftmargin=0pt,labelwidth=0pt,labelsep=0pt]
\item[\texttt{derive\_coefficients.py}]
Exact Bernoulli/power-sum derivation of $L_r(c)$ and $b_r(c)$, with a selectable
fixed order. The default is five.
\item[\texttt{asymptotics.py}]
Forward logarithmic profiles, the explicit cubic inverse, numerical inversion
of a chosen truncated profile, and the first exponential inverse correction.
\item[\texttt{verify.py}]
Exact product and divisor-recurrence checks, finite bounds, OEIS initial-value
checks, probability diagnostics, and inverse diagnostics.
\item[Data files]
\texttt{coefficients.json}, \texttt{exact\_values.json},\\
\texttt{critical\_errors.csv}, \texttt{poisson\_checks.csv},\\
\texttt{inverse\_checks.csv}, \texttt{verification\_summary.json}.
\end{description}
Only SymPy and mpmath are required beyond the Python standard library. The
article uses standard \LaTeX{} packages and builds with three runs of
\texttt{pdflatex}. The README supplies the exact commands. No network access
is required to reproduce the included calculations after dependencies are
installed.

\begin{writenote}[Added 1 October 2026, batch~74: the shipped layout]
In the collection the four programs (with \texttt{build.sh}) are in
\texttt{code/} and the data files, \texttt{requirements.txt} and the two
run records in \texttt{data/}; this source is \texttt{article.tex}, renamed
from \texttt{restricted\_partitions\_cubic\_boundary.tex}, and the
delivered PDF is not shipped. Every program reads and writes next to itself,
so \texttt{verify.py} cannot find \texttt{coefficients.json} from
\texttt{code/}, and \texttt{build.sh} still names the old source file.
Rerun on a flat copy of \texttt{code/} and \texttt{data/}; the README gives
the commands.
\end{writenote}

\section{A formalization route from the existing Lean foundation}

The exact combinatorial foundation already present in ProveIt suggests a
clean separation into finite and analytic layers. This is a proposed route,
not a claim of completed code.

The first layer would formalize the weak-composition orbit model and the
transposition fixed-point injection. These provide an entirely discrete proof
of the lower obstruction $m(m-1)^2/(4N)$. A companion lattice-simplex volume
lemma gives the upper bound and the exact threshold. The existing
multiplicity-vector equivalence can supply the bridge between the two models.

The second layer would formalize coefficient extraction for the finite
$q$-Pochhammer product, the analytic identity \eqref{rpc:eq:factor}, and the
Bernoulli recurrence \eqref{rpc:eq:a-rec}. At each fixed order, the explicit
coefficient identities reduce to rational polynomial algebra and are good
candidates for small proof certificates.

The genuinely analytic work is then localized: a uniform trigonometric
index-count lemma, a Cauchy-circle decay bound, gamma-kernel integral estimates,
and a uniform analytic Taylor remainder on a shrinking disk. The distinction
between a fixed parameter and a constant uniform over $N\ge c_0m^3$ must be
represented explicitly in the quantifiers. Finally, probability is built from
finite uniform measures on partition sets, coefficient-deletion identities,
and factorial moments.

The source file inspected in the repository is compatible with this program,
but neither its availability nor the successful Python checks certify any
of these future Lean developments. In particular, a formalization must not
replace the minor-arc estimate with the assertion that $q=1$ is the only
relevant singularity: the other roots of unity are present and must be bounded.

\section{Further research questions and concrete next steps}\label{rpc:sec:research}

The results above leave substantial, sharply formulated problems. The following
are proposals for further work, not assertions that every one is historically
open in exactly this form.

\paragraph{1. Optimal truncation and individual waves.}
Determine how the best order in \eqref{rpc:eq:centered} grows with $m$, and compute
the leading nontrivial root-of-unity contribution in the cubic regime. The
first target is the $q=-1$ wave, followed by a uniform comparison of all
remaining orders. This would turn the present fixed-order expansions into
a substantially more complete transseries.

\paragraph{2. A uniform bridge below the cubic scale.}
Extend the centered analysis into $m^{5/2}\lesssim N\ll m^3$, where the
centering factor can be large and the quadratic Bernoulli correction must
itself be exponentiated. Match the resulting expansion to a small-parameter
expansion of the Szekeres saddle. The precise validity range and matching
error are more useful targets than simply recovering a leading logarithm.

\paragraph{3. Structural degree cancellation in logarithmic coefficients.}
The general construction only immediately gives $\deg D_j\le3j$, while
$D_1$ and $D_2$ have degrees $3$ and $5$. Investigate a systematic sharper
degree theorem, its leading coefficients, and a connected-object or cumulant
interpretation. Such a result could reduce symbolic complexity and improve
error estimates across many dilute sequences.

\paragraph{4. Higher probability corrections and quantitative convergence.}
Compute an all-order version of \eqref{rpc:eq:pgf-correction}, then prove a
uniform complex-$u$ domain sufficient for extracting point probabilities
with explicit error bounds. Obtain a total-variation estimate for the Poisson
approximation, rather than inferring one from convergence of fixed moments.

\paragraph{5. Rare triple repetitions.}
The present proof bounds their probability by $O(1/m)$ at the cubic scale.
Determine its leading constant and joint asymptotics with the number of
doubletons. At other scales, identify when tripletons develop a nondegenerate
limit and whether a hierarchy of multiplicity transitions results.

\paragraph{6. Deletions whose number grows.}
Lemma~\ref{rpc:lem:delete} assumes fixed $k$. Find the maximal growing range
$k=k(m)$ for which a useful uniform coefficient ratio holds. This would give
moderate-deviation information for $Z_m$, not just a Poisson limit at fixed
counts.

\paragraph{7. Other weight systems.}
Replace $1,\ldots,m$ by structured integer weights with an available unit
weight and controlled power sums. The local expansion generalizes immediately,
but the global arc estimate depends on the weights' arithmetic distribution.
Find hypotheses that separate a universal volume/cumulant law from genuine
periodic obstructions.

\paragraph{8. Certified inverse thresholds.}
Combine explicit finite constants in the remainder bounds with monotonicity
to produce intervals enclosing the integer inverse $Q(y)$. A useful algorithm
must report when its interval straddles two integers, rather than silently
rounding a smooth chart. This is particularly important when $y$ lies very
close to an exact sequence value.

\paragraph{9. A proof-producing asymptotic pipeline.}
Starting from the finite-product generating function, automate power sums,
Bernoulli coefficients, fixed-order remainder parameters, transseries
substitution, and inverse reversion. The symbolic output should be a small
certificate checkable independently, with numerical fitting excluded from
the proof path.

\paragraph{10. OEIS and literature reconciliation.}
Prepare a narrowly scoped, source-backed update for the affected entries:
replace the vague quartic-growth condition by a precise sufficient condition,
attach the classical attribution, and separately record explicit refinements.
A useful contribution is improved mathematical traceability, not relabeling
classical consequences as breakthroughs.

\section{Conclusion}

The restricted-partition family admits one coherent explanation of several
OEIS asymptotic patterns. The volume approximation has an exact cubic
threshold, critical growth leaves a nontrivial constant, and larger sizes
admit polynomial, exponential, or factorial-scale refinements governed by the
same Bernoulli coefficient system. Uniform contour control is what turns
formal local algebra into theorems about the whole counting sequence.

The inverse formulas illustrate a second important distinction: an exact
gamma background is needed before exponentially small corrections can be
meaningfully inverted. Finally, the conjugate partition model shows that the
cubic boundary is not merely an artifact of asymptotic notation. It is a
regime where repeated parts survive in a nondegenerate Poisson law, with a
computable finite-size correction. The elementary criterion, all-order
expansions, inverse charts, and probabilistic law form the substantive
proof package supplied by this report.

\appendix
\section{A compact all-order coefficient recipe}

For a requested order $K$, the following is a finite exact procedure.
Compute $S_{2r}(m)$ for $1\le r\le K$ by Faulhaber's formula or finite
symbolic summation. Form $\gamma_r$, then $a_j$ from \eqref{rpc:eq:a-rec}.
Let $x=1/m$, $h=(x+x^2)/(4c)$, and form, modulo $x^{K+1}$,
\[
 Q(x,c)=\sum_{j=0}^K a_j(1/x)
 \frac{x^{4j}}{c^{2j}}\prod_{i=1}^{2j}(1-ix)\,(1+h)^{-2j}.
\]
Although the expression uses $a_j(1/x)$, the combined summand has no negative
powers: its lowest possible degree is $j$. Then
\begin{align*}
 \sum_{r=1}^K L_r(c)x^r\equiv{}&\log Q(x,c)
 +(x^{-1}-1)\log(1+h)-\frac1{4c}\\
 &-2\sum_{r\ge1}\frac{B_{2r}}{2r(2r-1)}x^{2r-1}
                          \pmod{x^{K+1}}.
\end{align*}
Only finitely many Stirling terms contribute. Apply \eqref{rpc:eq:b-rec} to obtain
relative coefficients. This is the algorithm implemented in
\texttt{derive\_coefficients.py}; it uses exact rational arithmetic.
For dilute coefficients, use \eqref{rpc:eq:h-rec}, \eqref{rpc:eq:beta}, and
\eqref{rpc:eq:D-rec} instead. No numerical sequence values are needed to discover
any of the displayed coefficients.

\begingroup
\small
\begin{thebibliography}{99}
\setlength{\itemsep}{3pt}
\bibitem{oeis238016}
OEIS Foundation Inc., \emph{A238016: Number of partitions of $n^k$ into parts
that are at most $n$}. Entry by Alois P.~Heinz; asymptotic comments by
V.~Kote\v{s}ovec (2015). Retrieved 1 October 2026.
\url{https://oeis.org/A238016}.

\bibitem{oeis238608}
OEIS Foundation Inc., \emph{A238608: Number of partitions of $n^3$ into parts
that are at most $n$}. Entry by Alois P.~Heinz; leading asymptotic and its
$c n^3$ extension by V.~Kote\v{s}ovec (2015). Retrieved 1 October 2026.
\url{https://oeis.org/A238608}.

\bibitem{oeis238010}
OEIS Foundation Inc., \emph{A238010: Number of partitions of $k^n$ into parts
that are at most $n$}. Entry by Alois P.~Heinz; asymptotic comment by
V.~Kote\v{s}ovec (2015). Retrieved 1 October 2026.
\url{https://oeis.org/A238010}.

\bibitem{oeis258670}
OEIS Foundation Inc., \emph{A258670: Number of partitions of $(2n)!$ into
parts that are at most $n$}. Retrieved 1 October 2026.
\url{https://oeis.org/A258670}.

\bibitem{canfield}
E.~Rodney Canfield, \emph{From recursions to asymptotics: On Szekeres' formula
for the number of partitions}, Electronic Journal of Combinatorics
\textbf{4}, no.~2 (1997), Research Paper 6, 16 pp.
\url{https://doi.org/10.37236/1321}.

\bibitem{dilchervignat}
Karl Dilcher and Christophe Vignat, \emph{An explicit form of the polynomial
part of a restricted partition function}, Research in Number Theory
\textbf{3} (2017), article 1.
\url{https://doi.org/10.1007/s40993-016-0065-3}.

\bibitem{sillszeilberger}
Andrew V.~Sills and Doron Zeilberger, \emph{Formulae for the number of
partitions of $n$ into at most $m$ parts (using the quasi-polynomial ansatz)},
arXiv:1108.4391 (2011).
\url{https://arxiv.org/abs/1108.4391}.

\bibitem{osullivan}
Cormac O'Sullivan, \emph{Partitions and Sylvester waves},
arXiv:1702.03611 (2017).
\url{https://arxiv.org/abs/1702.03611}.

\bibitem{dlmf}
NIST Digital Library of Mathematical Functions, \S5.11,
\emph{Asymptotic expansions of the gamma function}, especially 5.11.1
and the remainder bounds in \S5.11(ii). Retrieved 1 October 2026.
\url{https://dlmf.nist.gov/5.11}.

\bibitem{rpc-tai}
ProveIt, \emph{Transseries: the polynomial-logarithmic calculus, series
reversal at infinity, and inversion}, repository volume
\nolinkurl{Analysis/Transseries/docs/series-and-transseries/Transseries_And_Inversion/transseries_and_inversion.tex},
theorems labelled \texttt{p0:thm:staircase},
\texttt{p0:thm:perturbed-inversion} and \texttt{p6:thm:gamma}. Added in
writing (batch~74).

\bibitem{rpc-a097356}
ProveIt research-report collection,
\nolinkurl{SetTheory/Cardinals/docs/reports/generating-functions-and-asymptotics/oeis-sequence-asymptotics/a097356-sqrt-restricted-partitions},
research question ``Uniform limits as the saddle moves to zero or
infinity''. Unrefereed and not formalized. Added in writing (batch~74).

\bibitem{proveit}
Vladimir Reshetnikov, \emph{ProveIt repository}, inspected module
\texttt{PartitionBoundedParts.lean} (full path in Section~1.3),
main branch, retrieved 1 October 2026.
\url{https://github.com/VladimirReshetnikov/ProveIt}.
\end{thebibliography}
\endgroup
\end{document}
